\chapter{Introduction}
\label{ch:introduction}

\section{Context and Motivation}

The rapid evolution of Advanced Driver Assistance Systems (ADAS) and autonomous
driving architectures has progressively shifted the modern automotive industry
toward sensor-driven autonomy. Production vehicles are increasingly equipped with
ADAS functionalities such as adaptive cruise control, autonomous emergency
braking, and lane-keeping assistance, which actively mitigate human error and
improve traffic safety~\cite{yurtsever2020}. These systems depend entirely on
robust, real-time perception pipelines capable of constructing an accurate
three-dimensional representation of the surrounding environment. Perception is
not one module among many in an autonomous system: it is the entry point through
which all subsequent reasoning, planning, and control decisions flow, and
failures at this stage propagate irreversibly downstream.

Despite significant advances in single-agent perception, the standalone paradigm
faces a physical constraint that no algorithmic improvement can resolve: a single
vehicle cannot perceive beyond its line of sight. Occlusions behind large
commercial vehicles, blind intersections, and vehicles emerging from concealed
driveways represent everyday driving scenarios that leave autonomous systems
unaware of hazards until they enter the immediate sensor field of the
ego-vehicle. Cooperative perception addresses this at its root: by exchanging
spatial data over V2V communication networks, connected vehicles collectively
extend their sensing horizon beyond what any individual agent could achieve,
resolving blind spots through complementary viewpoints rather than better onboard
processing.

Deploying such a system at scale, however, raises an immediate practical question
of sensing modality. High-definition Light
Detection and Ranging (LiDAR) scanners have dominated the academic
literature due to their direct 3D measurements, but their cost, power
consumption, and mechanical complexity remain significant barriers to mass-market
deployment. Passive stereo camera arrays offer a compelling alternative: metric
depth recovery through geometric triangulation, without active illumination.
Their hardware cost and form factor are compatible with large-scale fleet
integration, making them a practical sensing backbone for cooperative systems.

It is this combination --- cooperative perception built on passive stereo
sensing --- that this project investigates.

\section{Problem Statement}

Deploying cooperative perception on passive stereo cameras raises a fundamental
feasibility concern. Stereo depth estimation recovers metric distance through
geometric triangulation, and its localisation error does not scale linearly with
range but quadratically. An object detected at 50 metres carries substantially
higher positional uncertainty than the same object detected at 15
metres~\cite{huang2025}. This range-dependent degradation means that the quality
of a stereo-sourced 3D detection is not a fixed property of the detector but a
function of where the target is relative to the sensing vehicle.

In a cooperative setting, this property becomes a structural problem. When two
stereo-equipped vehicles observe the same target from different distances, the
detections they produce carry fundamentally different uncertainty profiles.
Merging them into a unified environmental model requires reconciling estimates
whose reliability varies not just between agents but across the operational range
of each agent individually. Whether the geometric complementarity of two
viewpoints is sufficient to produce a net improvement in localisation accuracy,
and under what distance conditions that improvement holds or collapses, is not
established in the literature. Existing cooperative perception benchmarks and
evaluation frameworks assume LiDAR as the primary sensing modality, leaving the
stereo operating regime without a direct empirical
reference~\cite{xu2023v2v4real}.

The central question this work addresses is whether a cooperative perception
system built exclusively on passive stereo cameras can produce a measurable and
characterisable improvement in 3D localisation over single-agent stereo
perception, and what the operating limits of such an approach are.

\section{Research Objectives}

To evaluate how object-level late fusion handles range-dependent stereo noise,
this research defines three primary engineering and analytical objectives:

\begin{enumerate}[label=\textbf{O\arabic*.},leftmargin=*,itemsep=4pt]
  \item \textbf{Development of a Single-Agent 3D Spatial Localisation Baseline.}
  Implement a stereo vision pipeline that pairs 2D object detection with passive
  stereo depth estimation to project image-plane detections into a
  three-dimensional metric space. The objective is to establish the per-agent
  detection quality and localisation accuracy that feed into the cooperative
  layer.

  \item \textbf{Design of a Decentralised Object-Level Fusion Layer.} Construct a
  modular Vehicle-to-Vehicle (V2V) late fusion architecture that operates
  independently of agent software backbones. This module ingests 3D position
  estimates from multiple collaborating nodes, resolves the geometric
  transformations required to align independent coordinate frames, and merges
  redundant detections into a unified scene representation.

  \item \textbf{Empirical Characterisation of Cooperative Localisation
  Dynamics.} Conduct systematic experiments within real-world benchmark
  sequences and simulated multi-vehicle environments to map the performance
  profile of the cooperative stack. The focus is to quantify how spatial
  localisation error, precision, and recall evolve with increasing target
  distance, and to analyse how late fusion logic behaves when aggregating depth
  estimates from classical and learned stereo engines across agents.
\end{enumerate}

\section{Thesis Structure}

The remainder of this document is organised as follows to address these
objectives:

\begin{itemize}[leftmargin=*,itemsep=3pt]
  \item \textbf{Chapter~\ref{ch:sota}: Literature Review and State-of-the-Art}
  surveys existing work in stereo depth estimation, 3D spatial localisation, and
  collaborative perception to establish the boundaries of current research.
  \item \textbf{Chapter~\ref{ch:methodology}: Methodology and Proposed
  Framework} details the end-to-end engineering architecture of the pipeline and
  the V-model development process adopted to structure the design,
  implementation, and validation phases of the system.
  \item \textbf{Chapter~\ref{ch:experimental}: Experimental Setup and Simulation
  Environment} outlines the dataset configurations, simulated CARLA scenarios, and
  specific evaluation metrics used to test the system.
  \item \textbf{Chapter 5: Results and Discussion} presents the empirical data,
  performance plots, and the precision--recall trade-offs discovered during
  testing.
  \item \textbf{Chapter 6: Conclusion and Future Work} summarises the core
  findings, addresses current system limitations, and outlines future research
  pathways.
\end{itemize}
